\section{Dendritic solidification}
\label{sec:dendrite}

The 2012 Dendrite model starts with the whole puzzle already divided into
cells. Seeds expand by reaction and diffusion until neighboring regions meet.
A fixed figure reflects the diffusing fields, so a cell can wrap around its
outline. Two sinusoidal perturbations, one coarse and one fine, ripple the
initial seams. Their wavelengths and amplitudes help determine where branches
form \cite{puzzle2012}.

\paragraph{Why a ripple grows.}
Let \(p(\mathbf{x},t)\) be a local phase value between zero and one and
\(T(\mathbf{x},t)\) the dimensionless temperature at position \(\mathbf{x}\)
and time \(t\). The authors print the isotropic two-phase model in the form
\cite[Eqs.~5--7]{puzzle2012}
\begin{equation}
\begin{aligned}
  \tau\,\partial_t p
    &= \epsilon^2\nabla^2p+p(1-p)\bigl(p-\tfrac12+m(T)\bigr),\\
  \partial_t T
    &= \nabla^2 T+K\,\partial_t p,
    & m(T)&=\frac{\alpha}{\pi}\arctan\!\bigl(\gamma(T-T_e)\bigr).
\end{aligned}
\label{eq:dendrite-thermal}
\end{equation}
Here \(\epsilon\) sets the width of the diffuse interface, \(\tau\) its
relaxation time, and \(K\) the latent-heat coupling. The constants
\(\alpha\) and \(\gamma\) set the strength and temperature sensitivity of the
bias \(m\). The authors use \(T_e=0\) for equilibrium. At a growing tip,
heat diffuses sideways faster than the seam advances. The flanks warm
while the tip continues into colder material, amplifying the initial ripple
(Fig.~\ref{fig:thermal}).

\begin{figure}[!ht]\centering
\plate{thermal-feedback}{0.76}
\caption{Thermal feedback at an advancing tip. A protrusion reaches colder
material ahead. Heat from solidification spreads sideways and slows growth on
the flanks. The signs mark reinforcement and inhibition.}
\label{fig:thermal}
\end{figure}

\textit{Reading the sign convention.} The draft labels $p=0$ solid and
$p=1$ liquid, yet its positive $K\,\partial_t p$ term supplies heat as $p$
increases. Its Fig.~2 describes heat release on solidification. Equation~\eqref{eq:dendrite-thermal}
preserves the printed signs. A numerical reconstruction must choose a
consistent phase and heat convention~\cite[\S3.1]{puzzle2012}.

\paragraph{How neighboring pieces share a seam.}
A single solid--liquid front has only two sides. At a junction of three
pieces, each region needs a distinct identity. The model assigns each
active label \(i\) a phase fraction \(p_i\) and its own temperature \(T_i\),
with \(\sum_i p_i=1\). The draft combines interactions between pairs
\cite[Eqs.~8--10]{puzzle2012}. In compact notation,
\begin{equation}
  m_{ij}=\tfrac12\bigl[m(T_i)+m(T_j)\bigr],
  \qquad
  R_i=\sum_{j\ne i}p_i p_j\bigl(p_i-p_j+m_{ij}\bigr).
\label{eq:dendrite-pairwise}
\end{equation}
Here \(R_i\) denotes the summed local reaction for label \(i\). The factor
\(p_i p_j\) localizes each term to places where two phases overlap. With the
printed symmetric bias, the two ordered reactions for a pair sum to
$2p_i p_j m_{ij}$. An implementation must therefore also account for the
stated constraint $\sum_i p_i=1$. For
production-scale layouts the authors reused five labels across nonadjacent
pieces. Heat can then travel between pieces sharing a label and soften fine
branches in long runs. After growth, the studio traced the phase grid into
cut paths and smoothed tight folds to protect narrow wooden necks
\cite{puzzle2012,puzzleproject}.
